% GENERATED FILE. Edit canonical lessons, then run npm run generate:books.
% canonical-source-hash: fnv1a64:57552c8a33bd3d14
% canonical-lessons: MR-C541-savadhgiri, MR-C541-khabardari, MR-C541-daksata, MR-C541-satark, MR-C541-isara-dene

\chapter{Taking Care}
\label{ch:mr-taking-care}
\hlchaptermodality{541}

\begin{chapteropening}
\textbf{By the end of this chapter:} \emph{I can talk about caution, precautions and warnings.}

Complete the last lesson of chapter 541: I can talk about caution, precautions and warnings.
\end{chapteropening}

\section[\texorpdfstring{\mr{सावधगिरी} (sāvadhgirī)}{सावधगिरी (sāvadhgirī)}]{\mr{सावधगिरी} (sāvadhgirī) — caution}

\label{lesson:MR-C541-savadhgiri}



\begin{quote}
Before the new one: say the Marathi for an instalment, then the Marathi for to be affordable.
\end{quote}

\subsection*{You'll want to know: \mr{सावधगिरी}}
\textbf{\mr{सावधगिरी}} — \emph{sāvadhgirī} — \textquotedblleft{}caution\textquotedblright{}.

\begin{quote}
\mr{पावसात} \mr{गाडी} \mr{चालवताना} \mr{सावधगिरी} \mr{हवी}.
\end{quote}

\begin{quote}
\emph{pāvsāt gāḍī cālavtānā sāvadhgirī havī.}
\end{quote}

\begin{quote}
You need caution when driving in the rain.
\end{quote}

Feminine.

\begin{cousinweb}
\textbf{\mr{सावधगिरी}} (\emph{sāvadhgirī}) is built on \mr{सावध}, \textquotedblleft{}careful\textquotedblright{}, from Sanskrit \emph{sāvadhāna}, \textquotedblleft{}attentive\textquotedblright{} (Monier-Williams), with the Persian ending \emph{-girī}.
\end{cousinweb}

\begin{grammarlens}[title={What you've built}]
One more word for saying what happens if.
\end{grammarlens}

\subsection*{Your turn}
\begin{itemize}
\raggedright
  \item \emph{Say it:} \emph{sāvadhgirī}
  \item \emph{Say it:} \emph{sāvadhgirī}, once more
  \item \emph{Say it:} \emph{pāvsāt gāḍī cālavtānā sāvadhgirī havī.} — You need caution when driving in the rain.
  \item \emph{Recall:} say the Marathi for an instalment, then the Marathi for to be affordable, then say \emph{sāvadhgirī} again
\end{itemize}

\begin{tcolorbox}[breakable,colback=teal!4,colframe=teal!35!black,title={Before you move on}]
What does \mr{सावधगिरी} mean? (\textquotedblleft{}Caution\textquotedblright{}.) Say it once more.
\end{tcolorbox}
\section[\texorpdfstring{\mr{खबरदारी} (khabardārī)}{खबरदारी (khabardārī)}]{\mr{खबरदारी} (khabardārī) — a precaution; care}

\label{lesson:MR-C541-khabardari}



\begin{quote}
Before the new one: say the Marathi for to be affordable, then the Marathi for caution.
\end{quote}

\subsection*{You'll want to know: \mr{खबरदारी}}
\textbf{\mr{खबरदारी}} — \emph{khabardārī} — \textquotedblleft{}a precaution; care\textquotedblright{}.

\begin{quote}
\mr{आपण} \mr{खबरदारी} \mr{घेतली} \mr{तर} \mr{अपघात} \mr{टळतात}.
\end{quote}

\begin{quote}
\emph{āpaṇ khabardārī ghetlī tar apghāt ṭaḷtāt.}
\end{quote}

\begin{quote}
If we take care, accidents are avoided.
\end{quote}

Feminine.

\mr{खबरदारी} \mr{घेणे}, to take precautions.

\begin{cousinweb}
\textbf{\mr{खबरदारी}} (\emph{khabardārī}) is a loan from Persian \emph{khabardārī}, \textquotedblleft{}alertness\textquotedblright{}, built on \emph{khabar}, \textquotedblleft{}news\textquotedblright{}, which goes back to Arabic.
\end{cousinweb}

\begin{grammarlens}[title={What you've built}]
One more word for saying what happens if.
\end{grammarlens}

\subsection*{Your turn}
\begin{itemize}
\raggedright
  \item \emph{Say it:} \emph{khabardārī}
  \item \emph{Say it:} \emph{khabardārī}, once more
  \item \emph{Say it:} \emph{āpaṇ khabardārī ghetlī tar apghāt ṭaḷtāt.} — If we take care, accidents are avoided.
  \item \emph{Recall:} say the Marathi for to be affordable, then the Marathi for caution, then say \emph{khabardārī} again
\end{itemize}

\begin{tcolorbox}[breakable,colback=teal!4,colframe=teal!35!black,title={Before you move on}]
What does \mr{खबरदारी} mean? (\textquotedblleft{}A precaution; care\textquotedblright{}.) Say it once more.
\end{tcolorbox}
\section[\texorpdfstring{\mr{दक्षता} (dakṣatā)}{दक्षता (dakṣatā)}]{\mr{दक्षता} (dakṣatā) — vigilance}

\label{lesson:MR-C541-daksata}



\begin{quote}
Before the new one: say the Marathi for caution, then the Marathi for a precaution.
\end{quote}

\subsection*{You'll want to know: \mr{दक्षता}}
\textbf{\mr{दक्षता}} — \emph{dakṣatā} — \textquotedblleft{}vigilance\textquotedblright{}.

\begin{quote}
\mr{सणांच्या} \mr{दिवसांत} \mr{पोलीस} \mr{विशेष} \mr{दक्षता} \mr{घेतात}.
\end{quote}

\begin{quote}
\emph{saṇāñcyā divsānt polīs viśeṣ dakṣatā ghetāt.}
\end{quote}

\begin{quote}
On festival days the police are especially vigilant.
\end{quote}

Feminine.

\mr{दक्षता} \mr{घेणे}, to stay vigilant.

\begin{cousinweb}
\textbf{\mr{दक्षता}} (\emph{dakṣatā}) is a tatsama, taken from Sanskrit \emph{dakṣatā}, \textquotedblleft{}dexterity, ability\textquotedblright{} (Monier-Williams).
\end{cousinweb}

\begin{grammarlens}[title={What you've built}]
One more word for saying what happens if.
\end{grammarlens}

\subsection*{Your turn}
\begin{itemize}
\raggedright
  \item \emph{Say it:} \emph{dakṣatā}
  \item \emph{Say it:} \emph{dakṣatā}, once more
  \item \emph{Say it:} \emph{saṇāñcyā divsānt polīs viśeṣ dakṣatā ghetāt.} — On festival days the police are especially vigilant.
  \item \emph{Recall:} say the Marathi for caution, then the Marathi for a precaution, then say \emph{dakṣatā} again
\end{itemize}

\begin{tcolorbox}[breakable,colback=teal!4,colframe=teal!35!black,title={Before you move on}]
What does \mr{दक्षता} mean? (\textquotedblleft{}Vigilance\textquotedblright{}.) Say it once more.
\end{tcolorbox}
\section[\texorpdfstring{\mr{सतर्क} (satark)}{सतर्क (satark)}]{\mr{सतर्क} (satark) — alert; watchful}

\label{lesson:MR-C541-satark}



\begin{quote}
Before the new one: say the Marathi for a precaution, then the Marathi for vigilance.
\end{quote}

\subsection*{You'll want to know: \mr{सतर्क}}
\textbf{\mr{सतर्क}} — \emph{satark} — \textquotedblleft{}alert; watchful\textquotedblright{}.

\begin{quote}
\mr{जंगलात} \mr{फिरताना} \mr{सतर्क} \mr{राहिलं} \mr{पाहिजे}.
\end{quote}

\begin{quote}
\emph{jaṅglāt phirtānā satark rāhilaṁ pāhije.}
\end{quote}

\begin{quote}
One must stay alert when walking in the forest.
\end{quote}

\begin{cousinweb}
\textbf{\mr{सतर्क}} (\emph{satark}) is a tatsama, from Sanskrit \emph{satarka}, built on \emph{tarka}, \textquotedblleft{}reasoning\textquotedblright{}; it is glossed both as skilled in reasoning and as cautious (Monier-Williams).
\end{cousinweb}

\begin{grammarlens}[title={What you've built}]
One more word for saying what happens if.
\end{grammarlens}

\subsection*{Your turn}
\begin{itemize}
\raggedright
  \item \emph{Say it:} \emph{satark}
  \item \emph{Say it:} \emph{satark}, once more
  \item \emph{Say it:} \emph{jaṅglāt phirtānā satark rāhilaṁ pāhije.} — One must stay alert when walking in the forest.
  \item \emph{Recall:} say the Marathi for a precaution, then the Marathi for vigilance, then say \emph{satark} again
\end{itemize}

\begin{tcolorbox}[breakable,colback=teal!4,colframe=teal!35!black,title={Before you move on}]
What does \mr{सतर्क} mean? (\textquotedblleft{}Alert; watchful\textquotedblright{}.) Say it once more.
\end{tcolorbox}
\section[\texorpdfstring{\mr{इशारा} \mr{देणे} (iśārā deṇe)}{इशारा देणे (iśārā deṇe)}]{\mr{इशारा} \mr{देणे} (iśārā deṇe) — to issue a warning}

\label{lesson:MR-C541-isara-dene}



\begin{quote}
Before the new one: say the Marathi for vigilance, then the Marathi for alert.
\end{quote}

\subsection*{You'll want to know: \mr{इशारा} \mr{देणे}}
\textbf{\mr{इशारा} \mr{देणे}} — \emph{iśārā deṇe} — \textquotedblleft{}to issue a warning\textquotedblright{}.

\begin{quote}
\mr{हवामान} \mr{विभागाने} \mr{मोठ्या} \mr{पावसाचा} \mr{इशारा} \mr{दिला}.
\end{quote}

\begin{quote}
\emph{havāmān vibhāgāne moṭhyā pāvsācā iśārā dilā.}
\end{quote}

\begin{quote}
The weather department warned of heavy rain.
\end{quote}

Past, with -\mr{ने} (he …): \mr{त्याने} \mr{इशारा} \mr{दिला}, \emph{tyāne iśārā dilā}.

The danger takes -\mr{चा}: \mr{पावसाचा} \mr{इशारा} \mr{देणे}.

\begin{cousinweb}
\textbf{\mr{इशारा} \mr{देणे}} (\emph{iśārā deṇe}) is \mr{इशारा}, a loan from Persian \emph{išāra}, \textquotedblleft{}a sign, a hint\textquotedblright{}, usually traced to Arabic, with \mr{देणे}.
\end{cousinweb}

\begin{grammarlens}[title={What you've built}]
One more word for saying what happens if.
\end{grammarlens}

\subsection*{Your turn}
\begin{itemize}
\raggedright
  \item \emph{Say it:} \emph{iśārā deṇe}
  \item \emph{Say it:} \emph{iśārā deṇe}, once more
  \item \emph{Say it:} \emph{havāmān vibhāgāne moṭhyā pāvsācā iśārā dilā.} — The weather department warned of heavy rain.
  \item \emph{Recall:} say the Marathi for vigilance, then the Marathi for alert, then say \emph{iśārā deṇe} again
\end{itemize}

\begin{tcolorbox}[breakable,colback=teal!4,colframe=teal!35!black,title={Before you move on}]
What does \mr{इशारा} \mr{देणे} mean? (\textquotedblleft{}To issue a warning\textquotedblright{}.) Say it once more.
\end{tcolorbox}
